% REVTeX guide: https://www.ctan.org/tex-archive/macros/latex/contrib/revtex
\documentclass[aps,physrev,graphicx,amsmath,amssymb,reprint,superscriptaddress,nofootinbib,draft]{revtex4-2} % Two-column 'reprint'
%\documentclass[aps,physrev,graphicx,amsmath,amssymb,linenumbers,preprint]{revtex4-2}
\usepackage{graphicx}
\usepackage{subfig}
\usepackage{mathtools}
\usepackage{dcolumn}% Align table columns on decimal point
\newcolumntype{d}[1]{D{.}{\cdot}{#1}}
\usepackage{bm}% bold math
\usepackage{xcolor,soul}
\usepackage{physics}
\usepackage{tikz}
\usetikzlibrary{decorations.markings}
% FONT OPTIONS; mainly to break monotony when doing rounds of edits
\usepackage{mathpazo} %JMF - Palatino + math font :^)
%\usepackage{mlmodern} % 'fixed' computer modern with thicker weights
\usepackage{siunitx}% JMF addition - once a physicist, always a physicist
% Since version 2.1 siunitx supports so called comparators such as <, =, > and
% also ~= \approx and ~ \sim. Use as: \SI{\sim 200}{\kelvin}
%\bibliographystyle{aipnum4-1}
\newcommand{\nb}[1]{\textcolor{purple}{#1}}
\begin{document}
\title{Triplet Exciton Transfer: The Counterpoise Approach}
\date{\today} % useful to track hard copy of drafts

\author{Tom Ward}
\affiliation{Department of Chemistry, Imperial College London, Exhibition Road, London  SW7 2AZ, UK}
\author{Jarvist Moore Frost}
\affiliation{Department of Chemistry, Imperial College London, Exhibition Road, London  SW7 2AZ, UK}
\affiliation{Department of Physics, Imperial College London, Exhibition Road, London  SW7 2AZ, UK}
\email[Electronic mail:]{jarvist.frost@imperial.ac.uk}
\keywords{path-integral, Monte Carlo}
\begin{abstract}
\end{abstract}
\maketitle
\section{Introduction}\label{introduction}
\end{document}